\documentclass[11pt,a4paper]{article}
\usepackage{xeCJK}
\usepackage{fontspec}
\usepackage{xcolor}
\usepackage{geometry}
\usepackage[protrusion=true]{microtype}
\geometry{margin=1.8cm,top=1.9cm}
\linespread{1.25}

\setmainfont{Baskerville}
\setCJKmainfont{Songti SC}

\definecolor{tipblue}{RGB}{0,95,160}
\definecolor{rulegray}{RGB}{190,190,190}

\setlength{\parindent}{0pt}
\setlength{\parskip}{0.35em}
% 极长的英文单句（如第 135 期第 3 条）偶尔会差 <1pt 排不进一行、
% 且落在 \hangindent 悬挂段落里（\emergencystretch 对悬挂缩进无效）。
% \hfuzz 只抑制 <1pt 的、视觉不可见的水平溢出警告，不改变任何排版。
\hfuzz=1pt
\newcommand{\num}[1]{\makebox[1.9em][r]{#1.}\hspace{0.3em}}
\newcommand{\wline}{\noindent\rule{\linewidth}{0.4pt}}
% 中文题干：宋体粗体；提示词：蓝色小字；英文句子：Baskerville 正体
\newcommand{\prompt}[1]{\emph{\textcolor{tipblue}{\small（#1）}}}

\begin{document}
\pagestyle{plain}
\begin{center}
{\LARGE\bfseries 上海高考高分翻译背诵 \quad 134}\\[0.25em]
{\large 25二模　2026-11-24 \quad\ 默写（中 $\to$ 英）}\\[0.35em]
{\footnotesize 来源：微信公众号「发现之旅 速来记单词」· 2026-11-24 · 赵文瀚}
\end{center}

\vspace{0.4em}
\noindent\begin{tabular}{@{}p{0.33\textwidth}@{}p{0.34\textwidth}@{}p{0.33\textwidth}@{}}
姓名：\underline{\hspace{2.0cm}} & 日期：\underline{\hspace{2.0cm}} & 得分：\underline{\hspace{1.4cm}}\,/\enspace 4 \\
\end{tabular}

\vspace{0.5em}
\noindent{\small 根据中文句子与提示词写出英文翻译，题号沿用原文。}

\vspace{0.8em}
\num{1}\textbf{不要因为一次失败的尝试,就放弃你的目标。}\prompt{give}

\vspace{0.3em}

\wline\vspace{0.3em}
\wline

\num{2}\textbf{这位知名作家推荐的书广受追捧,现已售罄。}\prompt{so}

\vspace{0.3em}

\wline\vspace{0.3em}
\wline

\num{3}\textbf{这座北方冰城面对突如其来的大批游客,一时无所适从。}\prompt{number}

\vspace{0.3em}

\wline\vspace{0.3em}
\wline

\num{4}\textbf{当你知道这些交通工具和旅行线路背后的故事,你一定会惊叹人类一直以来是多么富有冒险精神。}\prompt{amazed}

\vspace{0.3em}

\wline\vspace{0.3em}
\wline

\end{document}
